\section{RELATED WORK AND PROBLEM FORMULATION}

\subsection{Rotating Legs and Wheeled Legs}

RHex showed that one continuously rotating compliant leg per hip can combine
mechanical simplicity with rugged mobility under a clock-driven tripod gait
\cite{saranli2001rhex,altendorfer2001rhex}; a dedicated controller later
demonstrated repeated ascent of full-size stairs \cite{moore2002reliable}.
The X-RHex family carried that morphology into a research platform
\cite{galloway2010xrhex}. Whegs abstracted
compliant, multi-spoke appendages and coupled propulsion to reduce actuator
count \cite{quinn2002whegs}; Q-Whex later used quasi-wheels and reported stair
climbing with geometry-dependent phase patterns \cite{zhang2023qwhex}. The
motion and power demand of C-legged hexapods have also been modeled and checked
on a physical test bench \cite{vina2021cleg}. The
half-circle RHex leg is the closest prior morphology to ours and the distinction
must be stated precisely: an RHex leg is a compliant C driven by a single hip
rotation, so its contact point and the body pose are determined by the same
input. Our C is the third joint of a 3-DoF leg, so the leg can hold the arc
still, place it, or spin it while the body is posed independently, and it is
this separation---not the C shape---that the ablation of Sec.~\ref{sec:results}
tests.

A second family transforms the mechanism itself. Wheel Transformer deploys
passive spokes when a wheel meets an obstacle \cite{kim2014wheeltransformer},
while Quattroped and TurboQuad change a leg into a wheel with a dedicated
transformation degree of freedom \cite{chen2014quattroped,chen2017turboquad}.
These recover both regimes at the cost of a transformation mechanism and its
actuation. A third family keeps a wheel drive alongside a full leg: ASTERISK H
switches a hexapod continuously between rolling and tripod motion using sensed
obstacles \cite{yoshioka2008asterisk}, while Cassino Hexapod III places
omni-wheels at anthropomorphic leg tips
\cite{tedeschi2017cassino}, ANYmal on Wheels integrates wheel constraints into
whole-body optimization \cite{bjelonic2019keep,bjelonic2021wholebody,hutter2016anymal},
and two-wheeled jumpers exploit the same separation in a smaller footprint
\cite{klemm2019ascento}. They demonstrate the value of simultaneous wheel and leg
control, but retain a wheel drive in addition to the leg chain. Our design does
not claim that hybrid locomotion is new. The distinction tested here is
\emph{actuator and contact reuse}: the distal leg joint itself produces the foot
trajectory, the rolling motion, and the edge engagement of one incomplete arc,
with no transformation state and no additional drive.

\subsection{Operating-Envelope Problem}

Let the body command be
\begin{equation}
  \vect{c}=[v_x,\;v_y,\;\omega_z]^\top\in\mathcal{C},
  \qquad m\in\{W,R,C\},                                      \label{eq:command}
\end{equation}
for walking, rolling, and climbing. Let $\vect{q}\in\R^{18}$ be the actuated
joint vector and $\Gamma_i$ the active contact subset of distal arc $i$. The
design seeks one morphology $\mathcal{M}$ and mode-dependent control laws
$\pi_m$ rather than a transformed mechanism $\mathcal{M}_m$:
\begin{equation}
 \vect{q}^{\star}_{k}=\pi_m(\vect{x}_k,\vect{c}_k;\mathcal{M}),
 \quad \mathcal{M}_W=\mathcal{M}_R=\mathcal{M}_C.             \label{eq:reuse}
\end{equation}
The relevant outcome is an operating envelope, not a single maximum speed,
\begin{equation}
 \mathcal{O}=\{\bar v_x, e_v, C_{\rm mt},D_{\rm slip},
 P_{\rm top}(h,d),t_{\rm top},t_{\rm switch}\},              \label{eq:envelope}
\end{equation}
where $h,d$ are riser and tread dimensions. We test three hypotheses.
H1: coordinated rolling reaches a higher mean forward speed than either walking
or distal-only rotation on the same model. H1 is deliberately stated as speed and
not as command tracking, because a controller that overshoots a fixed command
scores worse on $e_v$ while travelling faster; both quantities are reported
separately in Sec.~\ref{sec:results}. H2: on the open-arc model, coordinated
articulation reaches a stair profile on which distal-only rotation fails. A
claim about the opening itself requires the matched closed-wheel ablation
defined in Sec.~\ref{sec:method}.
H3: mode switching does not violate a $60^\circ$ tilt bound.

We also define the qualitative column of Table~\ref{tab:related}. Let $J_i$ map
leg $i$'s joint rates to the pair (tangential advance of the contact patch along
the distal surface, position of that contact relative to the body). Authority is
\emph{coupled} when $\operatorname{rank}J_i=1$, so that one input fixes both;
\emph{separated} when a dedicated drive makes the two independent at the cost of
an extra actuator; and \emph{partly separated} when, as here, three leg joints
can choose the pair over a bounded set but the distal joint is shared between
them.

\begin{table}[t]
\caption{Morphological position of the proposed platform}
\label{tab:related}
\centering\small
\resizebox{\columnwidth}{!}{%
\begin{tabular}{lcccc}
\toprule
Platform & leg actuation & rolling element & extra wheel drive & body/contact authority\\
\midrule
RHex \cite{saranli2001rhex} & 1/leg & compliant arc & no & coupled\\
Whegs I \cite{quinn2002whegs} & coupled & 3-spoke wheg & shared & coupled\\
Q-Whex \cite{zhang2023qwhex} & 1/leg & quasi-wheel & no & coupled\\
Cassino III \cite{tedeschi2017cassino} & articulated & omni-wheel & yes & separated\\
Proposed & 3/leg & open distal arc & no & partly separated\\
\bottomrule
\end{tabular}}
\end{table}
